% !TEX encoding = UTF-8 Unicode

% This is a simple template for a LaTeX document using the "article" class.
% See "book", "report", "letter" for other types of document.
\documentclass[12pt,oneside]{article}
 % use larger type; default would be 10pt


\usepackage[utf8]{inputenc} % set input encoding (not needed with XeLaTeX)

%%% Examples of Article customizations
% These packages are optional, depending whether you want the features they provide.
% See the LaTeX Companion or other references for full information.

%%% PAGE DIMENSIONS
\usepackage{geometry}

 % to change the page dimensions
\geometry{a4paper} % or letterpaper (US) or a5paper or....
%\geometry{margins=2in} % for example, change the margins to 2 inches all round
%\geometry{landscape} % set up the page for landscape
%   read geometry.pdf for detailed page layout information

%\usepackage{graphicx} % support the \includegraphics command and options

% \usepackage[parfill]{parskip} % Activate to begin paragraphs with an empty line rather than an indent

%%% PACKAGES
\usepackage{booktabs} % for much better looking tables
\usepackage{array} % for better arrays (eg matrices) in maths
\usepackage{paralist} % very flexible & customisable lists (eg. enumerate/itemize, etc.)
\usepackage{verbatim} % adds environment for commenting out blocks of text & for better verbatim
\usepackage{subfig} % make it possible to include more than one captioned figure/table in a single float
% These packages are all incorporated in the memoir class to one degree or another...

%%% HEADERS & FOOTERS
%\usepackage{fancyhdr} % This should be set AFTER setting up the page geometry



%%% SECTION TITLE APPEARANCE

 % (See the fntguide.pdf for font help)
% (This matches ConTeXt defaults)

%%% ToC (table of contents) APPEARANCE
%\usepackage[nottoc,notlof,notlot]{tocbibind} % Put the bibliography in the ToC
%\usepackage[titles,subfigure]{tocloft} % Alter the style of the Table of Contents
\usepackage[encapsulated]{CJK} 


\usepackage{setspace}
\usepackage{amsfonts}
\usepackage{amssymb}
\usepackage{amsmath}
\usepackage{amstext}
\usepackage{amscd}
\usepackage{hyperref}
\usepackage{amsthm}

\newcommand{\e}[0]{\varepsilon}
\newcommand{\norm}[1]{\left|\left|#1\right|\right|}
\newcommand{\ra}[0]{\rightarrow}
\newcommand{\Ra}[0]{\Rightarrow}
\newcommand{\R}[0]{\mathbb{R}}
\newcommand{\Q}[0]{\mathbb{Q}}
\newcommand{\N}[0]{\mathbb{N}}
\newcommand{\D}[0]{\mathbb{D}}
\newcommand{\BF}[0]{\mathbb{F}}
\newcommand{\CN}[0]{\mathcal{N}}
\newcommand{\Z}[0]{\mathbb{Z}}
\newcommand{\C}[0]{\mathbb{C}}
\newcommand{\CE}[0]{\mathcal{E}}
\newcommand{\CM}[0]{\mathcal{M}}
\newcommand{\CA}[0]{\mathcal{A}}
\newcommand{\CP}[0]{\mathcal{P}}
\newcommand{\CT}[0]{\mathcal{T}}
\newcommand{\cond}[4]{\left \{ \begin{array}{ll} #1 & \mbox{#2} \\ #3 & \mbox{#4} \end{array} \right.}
\newcommand{\condd}[6]{\left \{ \begin{array}{ll} #1 & \mbox{#2} \\ #3 & \mbox{#4} \\ #5 & \mbox{#6}\end{array} \right.}
\newcommand{\conddd}[8]{\left \{ \begin{array}{ll} #1 & \mbox{#2} \\ #3 & \mbox{#4} \\ #5 & \mbox{#6} \\ #7 & \mbox{#8} \end{array} \right.}
\newcommand{\lams}[1]{\lambda^*(#1)}
\newcommand{\bs}[0]{\backslash}
\newcommand{\bM}[0]{\overline{\M}}
\newcommand{\bmu}[0]{\overline{\mu}}
\newcommand{\mus}[0]{\mu^*}
\newcommand{\sa}[0]{\sigma \mbox{-algebra}}
\newcommand{\set}[1]{\{ #1\}}
\newcommand{\CB}[0]{\mathcal{B}}
\newcommand{\CG}[0]{\mathcal{G}}
\newcommand{\CF}[0]{\mathcal{F}}
\newcommand{\CU}[0]{\mathcal{U}}
\newcommand{\CL}[0]{\mathcal{L}}
\newcommand{\FA}[0]{\mathfrak{A}}
\newcommand{\BT}[0]{\mathbb{T}}
\newcommand{\BI}[0]{\mathbb{I}}
\newcommand{\BE}[0]{\mathbb{E}}
\newcommand{\eK}[1]{\e\mbox{-Kronecker}(#1)}
\newcommand{\eKG}[0]{\e\mbox{-Kronecker}}
\newcommand{\KL}[2]{\mbox{{\renewcommand*\rmdefault{}\normalfont KL}}(#1 \ || \ #2)}
\newcommand{\JS}[2]{\mbox{{\renewcommand*\rmdefault{}\normalfont JS}}(#1 \ || \ #2)}

\newcommand{\IN}[1]{I_0(#1)}
\newcommand{\ING}[0]{I_0}
\newcommand{\st}[0]{such that }
\newcommand{\CS}[0]{\mathcal{S}}
\newcommand{\CC}[0]{\mathcal{C}}
\newcommand{\tn}[1]{\text{\normalfont #1}}



\newcommand{\bmp}[1]{\overline{M}^+(#1)}
\newcommand{\simplep}[2]{S_{#1}^+(#2)}
\newcommand{\simple}[2]{S_{#1}(#2)}
\newcommand{\sumf}[2]{\sum_{#1 = #2}^\infty}
\newcommand{\bigcapf}[2]{\bigcap_{#1 = #2}^\infty}
\newcommand{\bigcupf}[2]{\bigcup_{#1 = #2}^\infty}
\newcommand{\limf}[1]{\lim_{#1 \ra \infty}}
\newcommand{\seq}[2]{\set{#1_{#2}}_{#2=1}^\infty}
\newcommand{\sg}[1]{\sigma\langle#1\rangle}
\newcommand{\ag}[1]{\langle#1\rangle}
\newcommand{\g}[1]{\langle#1\rangle}
\newcommand{\inte}[0]{\mbox{int}}
\newcommand{\Cl}[0]{\mbox{cl}}
\newcommand{\bd}[0]{\mbox{bd}}
\newcommand{\inv}[1]{{#1}^{-1}}
\newcommand{\Aut}[0]{\mbox{{\renewcommand*\rmdefault{}\normalfont Aut}}}
\newcommand{\Gal}[0]{\mbox{{\renewcommand*\rmdefault{}\normalfont Gal}}}
\newcommand{\Bor}[0]{\mbox{{\renewcommand*\rmdefault{}\normalfont Bor}}}
\newcommand{\Meas}[0]{\mbox{{\renewcommand*\rmdefault{}\normalfont Meas}}}
\newcommand{\ball}[2]{\mbox{{\renewcommand*\rmdefault{}\normalfont B}}_{#1}({#2})}
\newcommand{\supp}[0]{\mbox{{\renewcommand*\rmdefault{}\normalfont supp}}}
\newcommand{\Trig}[0]{\mbox{{\renewcommand*\rmdefault{}\normalfont Trig}}}
\newcommand{\sign}[0]{\mbox{{\renewcommand*\rmdefault{}\normalfont sign}}}
\newcommand{\AP}[0]{\mbox{{\renewcommand*\rmdefault{}\normalfont AP}}}

\newcommand{\ali}[1]{\begin{align*}#1\end{align*}}
\newcommand{\zhongwen}[1]{\begin{CJK}{UTF8}{gbsn}#1\end{CJK}}
\renewcommand{\Bar}[1]{\overline{#1}}
\renewcommand{\Hat}[1]{\widehat{#1}}
\newcommand{\Ga}[0]{\Gamma}
\newcommand{\ga}[0]{\gamma}
\newcommand{\inner}[2]{\left\langle #1, #2 \right\rangle}
\renewcommand{\Tilde}[1]{\widetilde{#1}}


\theoremstyle{plain}
\newtheorem{thm}{Theorem}[subsection]
\newtheorem{lem}[thm]{Lemma}
\newtheorem*{prop}{Proposition}
\newtheorem*{cor}{Corollary}
\newtheorem*{defn}{Definition}
\newtheorem{exmp}[thm]{Example}


\theoremstyle{remark}
\newtheorem*{rem}{Remark}
\newtheorem*{claim}{Claim}
\newtheorem*{note}{Note}
\newtheorem*{sol}{Solution}
\newtheorem*{case}{Case}


%%% END Article customizations

%%% The "real" document content comes below...

\title{Generative Adversarial Networks}
\date{} % Activate to display a given date or no date (if empty),
         % otherwise the current date is printed 
\parindent=0pt 
\begin{document}
\onehalfspace

\sloppy
\maketitle

\begin{defn}
Let $(E, \CA)$ and $(\Omega, \CM)$ be two measurable spaces. A Markov Kernel $\kappa$ with source $(E, \CA)$ and target $(\Omega, \CM)$ is a fucntion $\kappa: (E, \CM) \ra [0, 1]$ satisfying:\\
(1) for all $x \in E$, $\kappa(x)$ defines a probability measure on $(\Omega, \CM)$;\\
(2) for all $B \in \CM$, $\kappa(B): E \ra [0, 1]$ is $\CA$-measurable.
\end{defn}


Let $(\Omega, \CM, \mu_{\text{ref}})$ be a probability space. A {\bf generator} is a probability measure $\mu_G$ on $(\Omega, \CM)$, which is to mimic $\mu_{\text{ref}}$. The {\bf discriminator} is a Markov Kernel $\mu_D$ with source $(\Omega, \CM)$ and target $([0, 1], \CB([0, 1]))$, where $\CB([0, 1])$ is the Borel $\sa$. In other words: $\mu_D: \Omega \ra \CP([0, 1])$, where $\CP([0, 1])$ is the set of probability measures on $([0, 1], \CB([0, 1]))$. The discriminator is to tell either $x$ is from $\mu_{\text{ref}}$ or $\mu_G$.

The objective function is
\ali{L(\mu_G, \mu_D) = \BE_{x \sim \mu_{\text{ref}}, \ y \sim \mu_D(x)} [\ln y] + \BE_{x \sim \mu_G, \ y \sim \mu_D(x)} [\ln(1-y)].}

In the first expectation, the random variable $y$ with values in $[0, 1]$ follows the distribution $\mu_D(x)$, which is the posterior given $x$, and we need $y$ concentrates around $1$. In the second expectation, $y$ follows the posterior that $x$ is from the generator, and we need $y$ to be concentrated around $0$.

Therefore the generator tries to minimize $L$ and the discriminator tries to maximize $L$.

\begin{prop}
Given a generator $\mu_G$, 
\ali{\max_{\mu_D} L(\mu_G, \mu_D) = \max_{\substack{D: \ \Omega \ra [0,1] \\ D \ \tn{measurable}}} \BE_{x \sim \mu_\tn{ref}} [\ln D(x)] + \BE_{x \sim \mu_G} [\ln(1-D(x))]}
\end{prop}
\begin{proof}
Fix a generator $\mu_G$. Since $x \ra \ln(x)$ is concave, Jensen's inequality implies 
\ali{
    L(\mu_G, \mu_D) \leq \BE_{x \sim \mu_\tn{ref}} [\ln \BE_{y \sim \mu_D(x)} [y]] + \BE_{x \sim \mu_G} [\ln(1 - \BE_{y \sim \mu_D(x)}[y])].
}
Define $D(x) := \BE_{y \sim \mu_D(x)} [y]$. Since $D(x)$ is a pointwise limit of measurable simple functions, $D$ is measurable. This proves that LHS $\leq$ RHS.

Conversely, suppose $D:\Omega \ra [0,1]$ is measurable. We can define $\mu_D(x) = \delta_{D(x)}$ and $\mu_D(B) = I_{D^{-1}(B)}$ is measurable on $\Omega$. This shows LHS $\geq$ RHS.
\end{proof}

We then let the discriminator be a measurable function $D: \Omega \ra [0, 1]$. In practice, the generator arises from a prior distribution $z \sim \mu_z$ on a space $\Omega_z$ and a measurable fucntion $G:\Omega_z \ra \Omega$ such that $\mu_G = G^* \mu_z$. We update the objective function
\ali{
    L(\mu_G, D) = \BE_{x \sim \mu_\tn{ref}} [\ln D(x)] + \BE_{x \sim \mu_G} [\ln(1-D(x))].
}

\begin{prop}
Given a generator $\mu_G$, denote 
\ali{
    D^* = \tn{argmax}_{\substack{D: \ \Omega \ra [0, 1] \\ D \ \tn{measurable}}} L(\mu_G, D).
}
Then $D^*(x) = \frac{d\mu_\tn{ref}}{d(\mu_\tn{ref} + \mu_G)}(x)$.
\end{prop}
\begin{proof}
Let $\mu = \mu_\tn{ref} + \mu_G$. Define $p_\tn{ref} = \frac{d\mu_\tn{ref}}{d\mu}$ and $p_g = \frac{d\mu_G}{d\mu}$ as the Radon-Nikodym derivatives. Note that $p_g(x) + p_\tn{ref}(x) = 1$.
\ali{
    L(\mu_G, D) = \int_\Omega p_\tn{ref}(x) \ln D(x) + p_g(x) \ln(1-D(x)) \ d\mu(x).
}
The proof is completed as for all $x$, $D(x) = p_\tn{ref}(x)$ maximizes the integrand by the Gibbs' inequality.
\end{proof}

In this way we optimize
\ali{
    \min_{\mu_G} \max_{D} L(\mu_G, D) = \min_{\mu_G} L(\mu_G, D^*_G),
}
where $D^*_G$ is given above.

\begin{prop}
$\min_{\mu_G} L(\mu_G, D^*_G)$ is achieved when $\mu_G = \mu_\tn{ref}$. In this case $D^*_G(x) = 1/2$ for all $x$.
\end{prop}
\begin{proof}
Let $\mu$, $p_\tn{ref}$ and $p_g$ be as above.
\ali{
    L(\mu_G, D^*_G) &= \int_\Omega p_\tn{ref}(x) \ln p_\tn{ref}(x) + p_g(x) \ln p_g(x) \ d\mu(x) \\
                    &= \KL{\mu_\tn{ref}}{\mu} + \KL{\mu_G}{\mu} \\
                    &= 2\JS{\mu_\tn{ref}}{\mu_G} - 2\ln 2.
}
Hence, the minimum is obtained at $\mu_G = \mu_\tn{ref}$.
\end{proof}

Next we optimize $\max_{\mu_D} \min_{\mu_G} L(\mu_G, \mu_D)$.

\begin{prop}
We have
\ali{
\max_{\mu_D} \min_{\mu_G} L(\mu_G, \mu_D) = \min_{\mu_G} \max_{\mu_D} L(\mu_G, \mu_D) = -2 \ln 2
}
\end{prop}
\begin{proof}
We have 
\ali{
    \max_{\mu_D} \min_{\mu_G} L(\mu_G, \mu_D) \geq \max_{\mu_D} \BE_{x \sim \mu_{\text{ref}}, \ y \sim \mu_D(x)} [\ln y] + \inf_x \BE_{y \sim \mu_D(x)} [\ln(1-y)].
}
Similar to a previous proposition, 
\ali{
    &\max_{\mu_D} \BE_{x \sim \mu_{\text{ref}}, \ y \sim \mu_D(x)} [\ln y] + \inf_x \BE_{y \sim \mu_D(x)} [\ln(1-y)] \\
= &\max_{\substack{D: \ \Omega \ra [0,1] \\ D \ \tn{measurable}}} \BE_{x \sim \mu_\tn{ref}} [\ln D(x)] + \inf_x \ln(1-D(x)).
}
If we put $D(x) = 1/2$ for all $x$, we get 
\ali{
\max_{\mu_D} \min_{\mu_G} L(\mu_G, \mu_D) \geq -2 \ln 2 = \min_{\mu_G} \max_{\mu_D} L(\mu_G, \mu_D).
}
\end{proof}

To summarize,
\begin{cor}
$\mu_G^* = \mu_\tn{ref}$ and $\mu_D^* = \delta_{1/2}$ achieve the Nash equilibra:
\ali{
    \mu_G^* &= \tn{argmin}_{\mu_G} \ L(\mu_G, \mu_D^*) \\
    \mu_D^* &= \tn{argmax}_{\mu_D} \ L(\mu_G^*, \mu_D).
}
\end{cor}
\end{document}